\section{Background: diffusion models}
\label{sec:background}

Diffusion models \citep{sohl2015deep} are latent variable generative models characterized by a forward and a reverse Markov process. The forward process $q(\bx_{1:T}|\bx_0) = \prod_{t=1}^T q(\bx_t|\bx_{t-1})$ corrupts the data $\bx_0 \sim q(\bx_0)$ into a sequence of increasingly noisy latent variables $\bx_{1:T} = \bx_1, \bx_2, ..., \bx_T$. The learned reverse Markov process $p_{\theta}(\bx_{0:T}) = p(\bx_{T})\prod_{t=1}^T p_{\theta}(\bx_{t-1}|\bx_{t})$ gradually denoises the latent variables towards the data distribution. For example, for continuous data, the forward process typically adds Gaussian noise, which the reverse process learns to remove.

In order to optimize the generative model $p_{\theta}(\bx_{0})$ to fit the data distribution $q(\bx_0)$, we typically optimize a variational upper bound on the negative log-likelihood:
\begin{align}
    L_{\mathrm{vb}} = \mathbb E_{q(\bx_0)}\bigg[&
       \underbrace{D_{\mathrm{KL}}[q(\bx_T | \bx_0) \vert\vert p(\bx_T)]}_{L_T}
    + \sum_{t=2}^T \underbrace{\mathbb E_{q(\bx_t|\bx_0)} \big[
        D_{\mathrm{KL}}[q(\bx_{t-1} | \bx_t, \bx_0) \vert\vert 
        p_{\theta}(\bx_{t-1}|\bx_t)]
        \big]}_{L_{t-1}} \nonumber \\[-0.5em]
      &\underbrace{- \mathbb E_{q(\bx_1|\bx_0)} [\log p_{\theta}(\bx_0|\bx_1)]}_{L_0}
    \bigg].
    \label{eq:negative_lower_bound}
\end{align}

When the number of time steps $T$ goes to infinity, both the forward process and the reverse process share the same functional form \citep{feller1949theory}, allowing the use of a learned reverse process from the same class of distributions as that of the forward process.     
Furthermore, for several choices of the forward process the distribution $q(\bx_t|\bx_{0})$ converges to a stationary distribution $\pi(\bx)$ in the limit $t \rightarrow \infty$ independent of the value of $\bx_0$.
When the number of time steps $T$ is large enough and
we choose $\pi(\bx)$ as the prior $p(\bx_T)$,
we can guarantee that the 
$L_T$ term in \eqref{eq:negative_lower_bound} will approach zero regardless of the data distribution $q(\bx_0)$.
(Alternatively, one can use a learned prior $p_\theta(\bx_T)$.)

While $q(\bx_t | \bx_{t-1})$ can in theory be arbitrary, efficient training of $p_\theta$ is possible when $q(\bx_t | \bx_{t-1})$:
\begin{enumerate}
    \item Permits efficient sampling of $\bx_t$ from $q(\bx_t | \bx_0)$ for an arbitrary time $t$, allowing us to randomly sample timesteps and optimize each $L_{t-1}$ term individually with stochastic gradient descent,
    \item Has a tractable expression for the forward process posterior $q(\bs x_{t-1} | \bs x_{t}, \bs x_0)$, which allows us to compute the KL divergences present in the $L_{t-1}$ term of \eqref{eq:negative_lower_bound}.
\end{enumerate}
The majority of recent work in continuous spaces \citep{ho2020denoising, song2021implicit, wavegrad, nichol2021improved} defines the forward and reverse distributions as $q(\bx_t|\bx_{t-1}) = \mathcal N\left(\bx_{t} | \sqrt{1-\beta_t}\bx_{t-1}, \beta_t\bs I\right)$ and $p_{\theta}(\bx_{t-1}|\bx_{t}) = \mathcal N\left(\bx_{t-1} | \bs\mu_{\theta}(\bx_{t}, t), \bs \Sigma_{\theta}(\bx_{t}, t)\right)$, respectively. 
The aforementioned properties hold in the case of these Gaussian diffusion models:
the forward process $q(\bx_t | \bx_0)$ converges to a stationary distribution, motivating the choice $p(\bx_T) = \mathcal N\left(\bx_{T} | \bs 0, \bs I\right)$, and both $q(\bx_t|\bx_{0})$ and $q(\bs x_{t-1} | \bs x_{t}, \bs x_0)$ are tractable Gaussian distributions for which the KL divergence can be computed analytically.